\documentclass[10pt,a4paper]{amsart}
\usepackage[margin=19mm]{geometry}
\usepackage{amsmath,amssymb,mathtools}
\usepackage[scr=rsfs,cal=euler]{mathalfa}
\usepackage{xfrac}
\usepackage[unicode,hidelinks,bookmarks=false]{hyperref}
\newcommand{\ie}{i.e.\ }
\newcommand{\cf}{cf.\ }
\newcommand{\resp}{respectively\ }
\newcommand{\Subsection}{\subsection}
\providecommand{\nobreakdash}{\nobreak\hbox{-}}
\setlength{\emergencystretch}{2em}
\multlinegap0pt
\usepackage{bm}
\usepackage{amssymb}
\newtheorem{lemma}{Lemma}
\newcommand{\T}{\mathbb{T}}
\newcommand{\Z}{\mathbb{Z}}
\newcommand{\floor}[1]{\left\lfloor#1\right\rfloor}
\newcommand{\norm}[1]{\left\|#1\right\|}
\title{\textbf{Proof of the Finiteness of the Chromatic Number of Two-Dimensional Lacunary Integer Distance Graphs}}
\author{Nabh Singh}
\date{Source: arXiv:2606.22539v1 (CC BY 4.0)}
\begin{document}
\maketitle

\noindent\emph{Source and licence.} \textbf{Proof of the Finiteness of the Chromatic Number of Two-Dimensional Lacunary Integer Distance Graphs}. Version arXiv:2606.22539v1 (CC BY 4.0), distributed by arXiv under the Creative Commons Attribution 4.0 International licence, http://creativecommons.org/licenses/by/4.0/, as recorded in the arXiv licence metadata for this exact version and shown on the abstract page https://arxiv.org/abs/2606.22539v1. Full licence notice: \texttt{licenses/arxiv-2606.22539-CC-BY-4.0.txt}.\par\smallskip

\noindent\textit{Editorial scope.} Counted unit: the article's lemma lemma (source0.tex lines 85--92). The article's complete original proof of it is delivered with it (source0.tex lines 94--239, one \texttt{proof} environment of 146 lines). No mathematics is written, repaired or re-derived: the statement and the proof are the article's own lines, in the article's own notation, and no retained line is substituted or re-flowed.  No further source-local context is needed: every cross-reference of every retained line resolves inside the two delivered spans.  Nothing was omitted from any retained span, so the package carries no omission record.  No retained line contains a citation, so the wrapper carries no bibliography and no bibliography entry.  No retained line contains a figure environment, an image-inclusion command or drawing code, so no image is delivered and the package declares no asset dependency.  Editorial formatting only, in addition: an AMS \texttt{amsart} preamble that restates the article's own declarations and macros used by the retained lines, namely the environments \texttt{lemma} on separate counters, together with the standard packages those lines need.  The retained environments print in this document as Lemma 1 in order of appearance, whereas the article prints the same items with the numbering of its own full document.  The complete native source of the article as distributed in the arXiv e-print for this exact version is bundled unchanged as \texttt{sources/203355/source0.tex}.
\begin{lemma}[Torus-bin colouring lemma]
Let $D\subset \Z^2\setminus\{\mathbf 0\}$ be a lonely vector set with gap $\delta>0$. Then
\begin{equation}
    \chi(G(\Z^2,D))
    \le
    \left(\floor{\frac{2}{\delta}}+1\right)^2.
\end{equation}
\end{lemma}

\begin{proof}
Let $\boldsymbol\alpha=(\alpha_1,\alpha_2)\in [0,1)^2$ satisfy
\[
\norm{\langle \mathbf d,\boldsymbol\alpha\rangle}_{\T}\ge \delta
\qquad
\forall \mathbf d\in D.
\]

Define the coordinate-wise torus map
\[
\Phi_{\boldsymbol\alpha}:\Z^2\to \T^2
\]
by
\begin{equation}
\Phi_{\boldsymbol\alpha}(\mathbf v)
=
(v_1\alpha_1 \bmod 1,\; v_2\alpha_2 \bmod 1).
\end{equation}

We use the $\ell_\infty$ metric on $\T^2$, defined by
\[
\rho_\infty(\mathbf x,\mathbf y)
=
\max_{i=1,2}\norm{x_i-y_i}_{\T}.
\]

Let $\mathbf u,\mathbf v\in \Z^2$ be adjacent vertices in $G(\Z^2,D)$. Then
\[
\mathbf u-\mathbf v=\pm \mathbf d
\]
for some $\mathbf d\in D$. Write
\[
a_i=(u_i-v_i)\alpha_i
\qquad
(i=1,2).
\]
Then
\[
\langle \mathbf u-\mathbf v,\boldsymbol\alpha\rangle
=
a_1+a_2.
\]

By the triangle inequality on $\T$,
\[
\norm{a_1+a_2}_{\T}
\le
\norm{a_1}_{\T}+\norm{a_2}_{\T}.
\]
Therefore
\[
\norm{a_1+a_2}_{\T}
\le
2\max\{\norm{a_1}_{\T},\norm{a_2}_{\T}\}.
\]
\begin{equation}
\implies \max\{\norm{a_1}_{\T},\norm{a_2}_{\T}\}
\ge
\frac12\norm{a_1+a_2}_{\T}.
\end{equation}

Since
\[
a_1+a_2
=
\langle \mathbf u-\mathbf v,\boldsymbol\alpha\rangle,
\]
we obtain
\begin{align}
\rho_\infty\!\left(
\Phi_{\boldsymbol\alpha}(\mathbf u),
\Phi_{\boldsymbol\alpha}(\mathbf v)
\right)
&=
\max_{i=1,2}\norm{(u_i-v_i)\alpha_i}_{\T} \\
&\ge
\frac12
\norm{\langle \mathbf u-\mathbf v,\boldsymbol\alpha\rangle}_{\T} \\
&=
\frac12
\norm{\langle \mathbf d,\boldsymbol\alpha\rangle}_{\T} \\
&\ge
\frac{\delta}{2}.
\end{align}

Thus adjacent vertices are separated by at least $\delta/2$ in the $\ell_\infty$ torus metric.

Now choose
\[
M=\floor{\frac{2}{\delta}}+1.
\]
Then,
\[
\frac1M<\frac{\delta}{2}.
\]
Partition $[0,1)^2$ into $M^2$ half-open square bins
\[
Q_{ij}
=
\left[\frac{i}{M},\frac{i+1}{M}\right)
\times
\left[\frac{j}{M},\frac{j+1}{M}\right),
\]
where $0\le i,j\le M-1$.

Assign one colour to each square bin. Namely, define
\[
c:\Z^2\to \{0,\dots,M-1\}^2
\]
by
\[
c(\mathbf v)=(i,j)
\quad\Longleftrightarrow\quad
\Phi_{\boldsymbol\alpha}(\mathbf v)\in Q_{ij}
\]

If two vertices $\mathbf u,\mathbf v$ receive the same colour, then their images lie in the same square bin. In other words,
\[
\rho_\infty\!\left(
\Phi_{\boldsymbol\alpha}(\mathbf u),
\Phi_{\boldsymbol\alpha}(\mathbf v)
\right)
<
\frac1M
<
\frac{\delta}{2}
\]
However, adjacent vertices satisfy
\[
\rho_\infty\!\left(
\Phi_{\boldsymbol\alpha}(\mathbf u),
\Phi_{\boldsymbol\alpha}(\mathbf v)
\right)
\ge
\frac{\delta}{2}
\]
which is a contradiction.
Thus $c$ is a proper colouring of $G(\Z^2,D)$ using $M^2$ colours, and
\[
\chi(G(\Z^2,D))
\le
M^2
=
\left(\floor{\frac{2}{\delta}}+1\right)^2
\]
\end{proof}

\end{document}
